%% kapa: exam=1 round=21 year=2010 period=2 form=A subject=appraisal_law subject_ko=부동산관계법규 seq=1 number=41
%% source: Q-Net 감정평가사 기출문제 (cst00309) artlSeq=5024181 file=extracted/1008/5024181/2010년 제21회 감정평가사 1차시험 기출문제/2교시.hwp
%% conversion: kordoc 4.24.51 -> tools/gen1.py (auto); review=proofread-gpt-sol; flags=none
\begin{question}{41}
국토의 계획 및 이용에 관한 법령상 광역도시계획에 관한 설명으로 옳지 않은 것은?
\begin{choices}
  \item 광역도시계획이란 국토해양부장관 또는 도지사가 지정한 광역계획권의 장기발전방향을 제시하는 계획을 말한다.
  \item 광역도시계획은 국토해양부장관, 시ㆍ도지사, 시장 또는 군수가 수립한다.
  \item 광역도시계획을 수립하거나 변경하려면 미리 공청회를 열어 주민과 관계 전문가 등으로부터 의견을 들어야 한다.
  \item 시장 또는 군수는 광역도시계획을 수립하거나 변경하려면 도지사의 승인을 받아야 한다.
  \item 국토해양부장관, 시ㆍ도지사, 시장 또는 군수는 5년마다 광역도시계획에 대하여 그 타당성 여부를 재검토하여 이를 정비하여야 한다.
\end{choices}
\end{question}
\answer{5}
